%使用XeLaTex或者LualaLaTex编译

\documentclass[UTF8,AutoFakeBold=1,AutoFakeSlant,zihao=-4]{HZAU}
\usepackage{graphicx} % 用于插入图片
\usepackage{setspace} % 用于设置行距
\usepackage{ctex} % 支持中文字体和字号设置
\usepackage{color}
\usepackage{xcolor}
\usepackage[absolute, overlay]{textpos}
\usepackage{tcolorbox} % 用于创建带边框的盒子
\usepackage{fancyhdr}
\tcbuselibrary{skins} % 加载 skins 库以获得更多样式选项
\usepackage{multirow}
\usepackage{algorithm}
\usepackage{algpseudocode}
\usepackage{booktabs}

% 使用 natbib + gbt7714-numerical.bst 实现方括号上标引用
\usepackage[sort&compress,numbers,super]{natbib}
\renewcommand*{\citenumfont}[1]{[#1]}  % 给引用数字添加方括号
\bibliographystyle{gbt7714-numerical}


\definecolor{lightgray}{rgb}{.9,.9,.9}
\definecolor{darkgray}{rgb}{.4,.4,.4}
\definecolor{purple}{rgb}{0.65, 0.12, 0.82}
\definecolor{hzaugreen}{HTML}{67b356} 

%%%%%%%%%%%%%%%%
\usepackage{bicaption} % 支持双标题
% 设置中英文标题分隔符（默认为换行）
\DeclareCaptionOption{separator}{\def\captionseparator{#1}}
\captionsetup[bi-first]{name=图} % 中文"图"代替"Figure"
\captionsetup[bi-second]{name=Figure}
% 设置表格中英文标题格式
\captionsetup[table][bi-first]{name=表} % 中文"表"代替"Table"
\captionsetup[table][bi-second]{name=Table}

% 让 figure 编号按 section 分级（格式：Figure 1-1, 1-2, ...）
\numberwithin{figure}{section}
\numberwithin{table}{section}
\numberwithin{equation}{section} 

%%%%%%%%%%%%%%%请在这里填写你的论文题目
\newcommand{\thesisTitleC}{多酶级联催化绿色合成芳樟醇的研究}

\newcommand{\thesisTitleE}{Development of Multi-Enzyme Cascade Catalysis for Green Synthesis of Linalool}

%====================================================
%%%%%%%%%%%%%%请在这里填写你的个人信息
%====================================================
\newcommand{\yourMajor}{生物工程}  %专业
\newcommand{\yourMajorEn}{BIOENGINEERING}  %专业
\newcommand{\yourName}{张三}     %名字
\newcommand{\yourNameEn}{ZHANGSAN}     %名字
\newcommand{\yourMentor}{李四}        %导师
\newcommand{\yourMentorEn}{LISI}        %导师
\newcommand{\Mentorjob}{教授}           %职称
\newcommand{\MentorjobEn}{PROFESSOR}           %职称
\newcommand{\studentID}{2021317220XX}   %学号
\newcommand{\yourDateC}{二〇二五年六月}
\newcommand{\yourDateEn}{JUNE，2025}
\newcommand{\yourClass}{生物工程2101}       
\newcommand{\yourDegree}{工学学士学位}
\newcommand{\signatureDate}{2025年6月14日}
%% 记得提交最终版本时找自己导师要签名、以及自己的签名，补上这里的签名！！！
% 签名图片设置
% 如需添加签名，请将下方空白框替换为实际签名图片路径，如：figure/author_signature.png
\newcommand{\authorSignature}{example-image-empty}  % 作者签名占位
\newcommand{\mentorSignature}{example-image-empty}  % 导师签名占位
%====================================================

\begin{document}

% 封面（自动生成）
\coverpage      %%填写全部信息以确保no warning

%如果你需要印双面,请注意留白
\newpage
\thispagestyle{empty}  % 设置空白页不显示页眉页脚
\mbox{}  % 生成一个空白页
\newpage

\whitecoverpage      %%填写全部信息以确保no warning




%如果你需要印双面,请注意留白
\newpage
\thispagestyle{empty}  % 设置空白页不显示页眉页脚
\mbox{}  % 生成一个空白页
\newpage

\begin{authorization}
%此处为原创性声明,不用填写.也不用删除
\end{authorization}

\clearpage  % 强制清空浮动体
\thispagestyle{empty}
\hbox{}     % 比 \mbox{} 更底层
\vspace*{\fill}
\newpage

% 中文摘要
\begin{abstract}

% 此处为中文摘要内容，请在此填写你的论文摘要。
% 摘要应简要说明研究背景、研究目的、研究方法、主要结果和结论。
% 摘要字数一般在300-500字左右。

芳樟醇是一种重要的单萜类香料化合物，广泛应用于食品、化妆品和医药工业。传统的芳樟醇生产主要依赖植物提取和化学合成，但存在原料受限、环境污染等问题。本研究旨在开发一种多酶级联催化体系，实现芳樟醇的绿色生物合成。

本研究通过筛选和优化关键酶基因，构建了包含香叶基焦磷酸合酶（GPPS）和芳樟醇合酶（LIS）的双酶级联反应体系。通过优化酶的表达条件、反应温度、pH值和底物浓度等参数，显著提高了芳樟醇的转化率。同时，采用固定化酶技术提高酶的稳定性和重复使用率，降低了生产成本。

实验结果表明，优化后的多酶级联体系在30°C、pH 7.5条件下反应12小时，芳樟醇的摩尔转化率达到85.3\%，产物纯度超过98\%。固定化酶经过10批次反应后仍保持75\%以上的初始活性。本研究为芳樟醇的绿色合成提供了一种高效、环保的生物技术路线，具有重要的工业应用价值。
\keywords{芳樟醇；多酶级联催化；绿色合成；固定化酶；生物催化}%填写中文关键词,使用全角分号隔开
\end{abstract}

\begin{abstractEN}
%% add enligh abstract to table of contents
\addcontentsline{toc}{section}{Abstract}
%% must add a breaked line below to ensure the indent 

% This is the English abstract. Please write your abstract here.
% The abstract should briefly describe the research background, objectives, 
% methods, main results, and conclusions.

Linalool is an important monoterpene fragrance compound widely used in food, cosmetics, and pharmaceutical industries. Traditional linalool production mainly relies on plant extraction and chemical synthesis, which face issues such as limited raw materials and environmental pollution. This study aims to develop a multi-enzyme cascade catalysis system for the green biosynthesis of linalool.

In this study, a dual-enzyme cascade reaction system containing geranyl pyrophosphate synthase (GPPS) and linalool synthase (LIS) was constructed by screening and optimizing key enzyme genes. The conversion rate of linalool was significantly improved by optimizing enzyme expression conditions, reaction temperature, pH value, and substrate concentration. Meanwhile, immobilized enzyme technology was employed to enhance enzyme stability and reusability, reducing production costs.

Experimental results showed that the optimized multi-enzyme cascade system achieved a molar conversion rate of 85.3\% and product purity exceeding 98\% under conditions of 30°C and pH 7.5 for 12 hours. The immobilized enzyme maintained over 75\% of its initial activity after 10 batches of reactions. This study provides an efficient and environmentally friendly biotechnological route for linalool synthesis, with significant industrial application value.
\keywordsEN{Linalool; Multi-enzyme Cascade Catalysis; Green Synthesis; Immobilized Enzyme; Biocatalysis}  % 英文关键词
\end{abstractEN}


% 目录（自动生成）,不用修改
\contentpage


%%% 设置页眉
%%
% 正文开始
\pagestyle{fancy}
% 正文从第一页开始计算页码
\setcounter{page}{1}
% 页眉和页脚（页码）的格式设定
% 奇数页页眉显示论文题目，偶数页页眉显示学校信息
\fancyhf{}
\fancyhead[CO]{\fontsize{10.5pt}{10.5pt}\selectfont{\thesisTitleC}}  % 奇数页页眉：论文题目
\fancyhead[CE]{\fontsize{10.5pt}{10.5pt}\selectfont{华中农业大学2025届学士学位毕业论文}}  % 偶数页页眉
\fancyfoot[C]{\fontsize{9pt}{9pt}\selectfont{\thepage}}
\renewcommand{\headrulewidth}{1pt}
\renewcommand{\footrulewidth}{0pt}

% 正文 22 磅的行距，段前段后间距为 0
\setlength{\parskip}{0em}
\renewcommand{\baselinestretch}{1.53}
% 正文首行悬挂 1.02cm
\setlength{\parindent}{1.02cm}
%%%%%%%%%%%%%%%%%%%以下是正文部分%%%%%%%%%%%%%%%%%%%%%%%

%% 注意，按照毕设要求，引用文章使用\citep而非cite
\section{引言}

\subsection{研究背景和意义}

% 请在此处填写研究背景和意义
% 介绍研究领域的背景、现状以及本研究的重要性和意义

本文研究了多酶级联催化绿色合成芳樟醇的问题。随着生物催化技术的发展，绿色合成方法越来越受到关注。本研究对于芳樟醇的可持续生产具有重要意义。

\subsection{国内外研究现状}

\subsubsection{芳樟醇生物合成研究现状}

% 请在此处介绍相关领域的研究现状
% 可以从不同角度或方法分类介绍现有研究工作

（1）\textbf{基于植物提取的方法}

现有研究在植物提取芳樟醇方面取得了一定进展\cite{dudareva2013biosynthesis}。这些方法主要关注从薰衣草等植物中提取芳樟醇，但在原料来源和成本方面仍存在不足。

（2）\textbf{基于化学合成的方法}

另一类方法采用化学合成策略\cite{peralta2012identification}，取得了较好的效果。然而，这些方法在环境污染和选择性方面仍有改进空间。

（3）\textbf{基于生物催化合成的方法}

近年来，生物催化法合成芳樟醇成为研究热点\cite{caron2019biosynthesis}，具有反应条件温和、选择性高等优点。

\subsubsection{研究现状总结}

% 总结现有研究的不足和本研究的切入点

综上所述，现有研究在芳樟醇合成方面存在以下不足：（1）植物提取法原料受限、成本高；（2）化学合成法环境污染严重；（3）生物催化法的转化效率和酶稳定性有待提高。针对这些问题，本文提出了多酶级联催化绿色合成芳樟醇的方法。

\subsection{研究内容}

% 请在此处列出本文的主要研究内容

针对上述问题，本文围绕多酶级联催化合成芳樟醇展开研究，主要研究内容包括：

\textbf{（1）关键酶的筛选与表达}

% 介绍第一个研究内容

筛选高活性的香叶基焦磷酸合酶（GPPS）和芳樟醇合酶（LIS），通过基因工程手段在大肠杆菌中实现高效异源表达。

\textbf{（2）多酶级联反应体系的构建与优化}

% 介绍第二个研究内容

构建GPPS和LIS的双酶级联反应体系，系统研究反应条件对芳樟醇转化率的影响。

\textbf{（3）固定化酶的制备与性能评价}

% 介绍第三个研究内容

采用合适的载体和固定化方法制备固定化酶，评价固定化酶的活性、稳定性、重复使用性能等关键指标。

\textbf{（4）多酶级联反应的放大与工艺评价}

% 介绍第四个研究内容（如有）

在优化的反应条件下进行反应放大试验，评价多酶级联催化体系的工业化应用潜力。

\subsection{论文结构}

% 请在此处介绍论文的组织结构

本文共分为五章，各章节内容安排如下：

第一章为引言，介绍研究背景和意义，综述国内外研究现状，明确本文的研究内容和论文组织结构。

第二章为相关理论与关键技术，介绍萜类化合物生物合成途径、酶工程技术、固定化酶技术等基础知识，为后续章节的方法设计奠定理论基础。

第三章为多酶级联催化体系的构建与优化，详细阐述关键酶的筛选与表达、级联反应体系的构建、反应条件优化等内容。

第四章为实验结果与分析，介绍实验材料与方法、酶的表达与纯化结果、级联反应优化结果、固定化酶性能评价等。

第五章为总结与展望，总结论文的主要工作，分析研究的局限性，并对未来研究方向进行展望。
 \clearpage

\section{相关理论与关键技术}

% 本章介绍与本文研究相关的理论基础与关键技术

本章介绍与本文研究相关的理论基础与关键技术，包括萜类化合物的生物合成途径、酶工程技术的基本原理以及固定化酶技术的方法与应用。

\subsection{萜类化合物的生物合成}

\subsubsection{萜类化合物的结构与分类}

% 请在此处介绍萜类化合物的基本知识

萜类化合物（Terpenoids）是一类以异戊二烯（C5H8）为基本结构单元的天然产物。芳樟醇属于单萜类化合物，分子式为C10H18O。

\subsubsection{萜类生物合成的两条主要途径}

萜类化合物的生物合成主要通过两条途径实现：甲羟戊酸（MVA）途径和甲基赤藓糖醇磷酸（MEP）途径。

\textbf{（1）甲羟戊酸（MVA）途径}

MVA途径主要存在于细胞质中，以乙酰辅酶A为起始底物。

\textbf{（2）甲基赤藓糖醇磷酸（MEP）途径}

MEP途径主要存在于质体中，以丙酮酸和甘油醛-3-磷酸为起始底物。

\subsubsection{单萜的生物合成}

单萜的生物合成以IPP和DMAPP为前体。首先，IPP和DMAPP在香叶基焦磷酸合酶（GPPS）催化下缩合生成香叶基焦磷酸（GPP，C10）。然后，芳樟醇合酶（LIS）催化GPP转化为芳樟醇。

该反应的化学方程式可表示为：

\begin{equation}
\text{GPP} \xrightarrow{\text{LIS}} \text{Linalool} + \text{PP}_i
\end{equation}

\subsection{酶工程技术}

\subsubsection{酶的异源表达}

% 请在此处介绍酶的异源表达技术

酶的异源表达是酶工程的核心技术之一，通过将目标酶基因克隆到表达载体中，转化到合适的宿主细胞中实现目标酶的过量表达。

\subsubsection{酶的纯化与活性测定}

% 请在此处介绍酶的纯化与活性测定方法

酶的纯化通常采用亲和层析、离子交换层析、凝胶过滤层析等方法。酶活力单位（U）定义为在特定条件下每分钟催化生成1 μmol产物所需的酶量。

\subsection{固定化酶技术}

\subsubsection{固定化酶的制备方法}

% 请在此处介绍固定化酶的制备方法

固定化酶技术是将酶固定在特定载体上，使酶在保持催化活性的同时便于回收和重复使用。常用的固定化方法包括：

\textbf{（1）吸附法}

吸附法是通过物理作用力将酶吸附在载体表面。

\textbf{（2）包埋法}

包埋法是将酶包埋在凝胶网格或微囊结构中。

\textbf{（3）共价结合法}

共价结合法是通过酶的氨基酸残基与载体表面的活性基团形成共价键将酶固定。

\textbf{（4）交联法}

交联法是使用双功能或多功能交联剂使酶分子之间或酶与载体之间形成共价交联。

\subsubsection{固定化酶的性能评价}

% 请在此处介绍固定化酶的性能评价指标

固定化酶的性能主要通过以下指标评价：固定化率、相对活力、操作稳定性、储存稳定性、热稳定性等。

\subsection{本章小结}

% 总结本章内容

本章介绍了与本文研究相关的关键技术和理论基础，包括萜类化合物的分类和生物合成途径、酶工程技术、固定化酶技术等。这些技术为后续章节构建多酶级联催化体系奠定了理论基础。
 \clearpage

\section{多酶级联催化体系的构建与优化}

% 本章详细介绍多酶级联催化绿色合成芳樟醇的方法

本章详细介绍多酶级联催化绿色合成芳樟醇的方法。

\subsection{概述}

% 请在此处给出方法的整体概述

本文提出的多酶级联催化体系如图 \ref{fig:framework} 所示，以IPP和DMAPP为底物，通过GPPS和LIS的双酶级联催化合成芳樟醇。

\subsection{关键酶的筛选与表达}

\subsubsection{酶基因的筛选}

% 请在此处介绍酶基因的筛选过程

GPPS和LIS是芳樟醇生物合成的两个关键酶。本研究从文献报道和数据库中筛选了来源于不同物种的GPPS和LIS基因序列。

\subsubsection{表达载体的构建}

% 请在此处介绍表达载体的构建方法

将筛选得到的基因序列进行密码子优化，连接到表达载体pET-28a(+)中，构建重组表达质粒。

\subsubsection{诱导表达与条件优化}

% 请在此处介绍诱导表达条件

将含有重组质粒的大肠杆菌进行诱导表达，通过SDS-PAGE和Western blot验证目标蛋白的表达。

\subsection{多酶级联反应体系的构建}

\subsubsection{级联反应体系的设计}

% 请在此处介绍级联反应体系的设计

双酶级联反应体系包含GPPS和LIS两种酶，以IPP和DMAPP为起始底物。

\subsubsection{反应条件优化}

% 请在此处介绍反应条件优化

系统研究了反应温度、pH值、底物浓度、酶比例、反应时间等因素对芳樟醇转化率的影响。

\textbf{（1）温度优化}

在20-50°C范围内设置不同温度梯度进行反应，确定最优反应温度。

\textbf{（2）pH值优化}

在pH 6.0-9.0范围内设置不同pH梯度进行反应，确定最优反应pH值。

\textbf{（3）底物浓度优化}

设置不同浓度的IPP和DMAPP进行反应，确定最优底物浓度。

\textbf{（4）酶比例优化}

设置不同的酶摩尔比，测定芳樟醇的生成速率，确定最优酶比例。

\subsection{固定化酶的制备}

\subsubsection{载体的选择}

% 请在此处介绍载体的选择

选择氨基功能化的磁性纳米粒子作为固定化载体。

\subsubsection{固定化方法}

% 请在此处介绍固定化方法

采用戊二醛交联法固定化酶。

\subsubsection{固定化条件的优化}

% 请在此处介绍固定化条件的优化

优化了载体用量、酶载量、固定化pH值和固定化时间等参数。

\subsection{固定化酶的性能评价}

\subsubsection{酶学性质测定}

% 请在此处介绍酶学性质测定

测定了固定化酶的最适温度、最适pH、热稳定性等酶学性质。

\subsubsection{操作稳定性评价}

% 请在此处介绍操作稳定性评价

在优化条件下，评价了固定化酶的重复使用性能。

\subsection{本章小结}

% 总结本章内容

本章详细介绍了多酶级联催化绿色合成芳樟醇的方法，包括关键酶的筛选与表达、级联反应体系的构建与优化、固定化酶的制备与性能评价。
 \clearpage

\section{实验结果与分析}

% 本章通过实验验证多酶级联催化体系合成芳樟醇的有效性

本章通过实验验证多酶级联催化体系合成芳樟醇的有效性。

\subsection{实验材料与方法}

\subsubsection{实验材料}

% 请在此处介绍实验材料

\textbf{（1）菌株与质粒}

大肠杆菌DH5α用于基因克隆，大肠杆菌BL21(DE3)用于蛋白表达。

\textbf{（2）主要试剂}

IPP、DMAPP、GPP标准品、芳樟醇标准品等。

\textbf{（3）主要仪器}

高效液相色谱仪、气相色谱-质谱联用仪、紫外-可见分光光度计等。

\subsubsection{分析方法}

% 请在此处介绍分析方法

\textbf{（1）芳樟醇的定量分析}

采用气相色谱-质谱联用（GC-MS）方法定量分析芳樟醇含量。

\textbf{（2）酶活性的测定}

GPPS和LIS的活性测定方法描述。

\subsection{酶的表达与纯化}

\subsubsection{重组酶的表达}

% 请在此处介绍重组酶的表达结果

SDS-PAGE分析结果显示，GPPS和LIS在IPTG诱导后均有明显表达。

\subsubsection{酶的纯化}

% 请在此处介绍酶的纯化结果

采用镍亲和层析纯化重组酶，纯化结果如表 \ref{tab:purification} 所示。

\begin{table}[h]
  \centering
  \bicaption{酶的纯化结果}{Purification results of enzymes}
  \label{tab:purification}
  \begin{tabular}{lccccc}
    \toprule
    \textbf{酶} & \textbf{纯化步骤} & \textbf{总蛋白(mg)} & \textbf{比活力(U/mg)} & \textbf{纯化倍数} & \textbf{回收率(\%)} \\
    \midrule
    GPPS & 粗酶液 & 120.5 & 2.8 & 1.0 & 100 \\
         & 纯化后 & 18.2 & 15.6 & 5.6 & 84.3 \\
    \midrule
    LIS & 粗酶液 & 95.3 & 1.5 & 1.0 & 100 \\
        & 纯化后 & 12.8 & 8.2 & 5.5 & 78.6 \\
    \bottomrule
  \end{tabular}
\end{table}

\subsection{级联反应条件优化}

\subsubsection{温度对反应的影响}

% 请在此处介绍温度优化结果

考察了20-50°C范围内温度对芳樟醇转化率的影响，确定最优反应温度。

\subsubsection{pH对反应的影响}

% 请在此处介绍pH优化结果

考察了pH 6.0-9.0范围内pH对芳樟醇转化率的影响，确定最优反应pH值。

\subsubsection{底物浓度对反应的影响}

% 请在此处介绍底物浓度优化结果

考察了不同底物浓度对芳樟醇转化率的影响，确定最优底物浓度。

\subsubsection{酶比例对反应的影响}

% 请在此处介绍酶比例优化结果

考察了GPPS和LIS不同摩尔比对反应的影响，结果如表 \ref{tab:enzyme_ratio} 所示。

\begin{table}[h]
  \centering
  \bicaption{酶比例对芳樟醇转化率的影响}{Effect of enzyme ratio on linalool conversion}
  \label{tab:enzyme_ratio}
  \begin{tabular}{ccc}
    \toprule
    \textbf{GPPS:LIS (摩尔比)} & \textbf{芳樟醇转化率(\%)} & \textbf{GPP积累量(mM)} \\
    \midrule
    1:1 & 72.5 & 0.18 \\
    1:2 & 78.3 & 0.12 \\
    1:5 & 82.6 & 0.05 \\
    1:10 & 83.1 & 0.02 \\
    2:1 & 65.8 & 0.35 \\
    5:1 & 58.2 & 0.52 \\
    10:1 & 52.5 & 0.68 \\
    \bottomrule
  \end{tabular}
\end{table}

\subsection{最优条件下的反应结果}

% 请在此处介绍最优条件下的反应结果

在优化条件下，芳樟醇的摩尔转化率达到85.3\%，产物纯度超过98\%。

\subsection{固定化酶的性能评价}

\subsubsection{固定化条件优化}

% 请在此处介绍固定化条件优化结果

优化了载体用量、酶载量、固定化pH值和固定化时间等参数。

\subsubsection{固定化酶的酶学性质}

% 请在此处介绍固定化酶的酶学性质

测定了固定化酶的最适温度、最适pH、热稳定性等酶学性质。

\subsubsection{固定化酶的操作稳定性}

% 请在此处介绍固定化酶的操作稳定性

评价了固定化酶的重复使用性能。

\subsection{本章小结}

% 总结本章内容

本章通过实验验证了多酶级联催化体系合成芳樟醇的有效性，确定了最优反应条件，制备的固定化酶具有良好的稳定性。
 \clearpage

\section{总结与展望}

\subsection{论文工作总结}

% 请在此处总结论文工作

本文针对芳樟醇传统生产方法存在的问题，开展了多酶级联催化绿色合成芳樟醇的研究。本文的主要工作总结如下：

\textbf{（1）成功筛选并表达了芳樟醇合成的关键酶}

% 总结第一个贡献

从薄荷和薰衣草中筛选了高活性的GPPS和LIS基因，通过基因工程手段在大肠杆菌中实现了高效异源表达。

\textbf{（2）构建了高效的多酶级联催化体系}

% 总结第二个贡献

构建了GPPS和LIS的双酶级联反应体系，通过系统优化反应条件，芳樟醇的摩尔转化率达到85.3\%。

\textbf{（3）开发了固定化酶技术提高酶的稳定性}

% 总结第三个贡献

采用磁性纳米粒子固定化技术制备了固定化酶，显著提高了酶的稳定性和重复使用性能。

\textbf{（4）建立了芳樟醇绿色合成的技术路线}

% 总结第四个贡献

本研究建立了一条从简单前体IPP和DMAPP出发，通过多酶级联催化合成芳樟醇的绿色技术路线。

\subsection{研究局限性和未来展望}

% 请在此处讨论局限性和未来方向

尽管本研究取得了一定的成果，但仍存在一些局限性。

\textbf{（1）当前研究的局限性}

\textit{局限性一：底物成本较高}。

% 描述第一个局限性

本研究使用的底物IPP和DMAPP价格较高，限制了该方法的经济可行性。

\textit{局限性二：酶活性有待进一步提高}。

% 描述第二个局限性

目前GPPS和LIS的催化效率仍有提升空间。

\textbf{（2）未来研究方向}

基于上述局限性，未来研究可以从以下几个方向展开：

\textit{方向一：构建全细胞催化体系}。

% 描述第一个未来方向

将GPPS和LIS基因整合到同一宿主菌中，构建产芳樟醇的工程菌株。

\textit{方向二：酶的定向进化与改造}。

% 描述第二个未来方向

利用定向进化、理性设计等蛋白质工程技术，对GPPS和LIS进行改造。

\textit{方向三：反应体系的放大与工艺优化}。

% 描述第三个未来方向

开展反应体系的放大研究，建立适合工业化生产的生产工艺。
 \clearpage

% 显示所有参考文献，即使未被引用,如果你只想引用时显示,请删除下面的 \nocite{*}
\nocite{*}
\bibliography{references}


%%%%%%附录
%% 根据个人需要是否添加appendix
\appendix       %开始附录部分
\phantomsection % 手动添加锚点，确保超链接正确

\section{主要符号说明}

% 本附录列出论文中使用的关键符号及其含义

\begin{table}[h]
  \centering
  \bicaption{主要符号说明}{Notation definitions}
  \label{tab:notation}
  \begin{tabular}{cl}
    \toprule
    \textbf{符号} & \textbf{含义} \\
    \midrule
    GPPS & 香叶基焦磷酸合酶 \\
    LIS & 芳樟醇合酶 \\
    IPP & 异戊烯基焦磷酸 \\
    DMAPP & 二甲基烯丙基焦磷酸 \\
    GPP & 香叶基焦磷酸 \\
    MVA & 甲羟戊酸途径 \\
    MEP & 甲基赤藓糖醇磷酸途径 \\
    IPTG & 异丙基-β-D-硫代半乳糖苷 \\
    $U$ & 酶活力单位 \\
    $\mu$mol & 微摩尔 \\
    \bottomrule
  \end{tabular}
\end{table}

\section{常用缓冲液配方}

\subsection{LB培养基}

\begin{itemize}
    \item 胰蛋白胨：10 g/L
    \item 酵母提取物：5 g/L
    \item NaCl：10 g/L
    \item pH 7.0，121°C高压灭菌20 min
\end{itemize}

\subsection{酶活性测定缓冲液}

\begin{itemize}
    \item Tris碱：6.06 g
    \item 双蒸水：800 mL
    \item 用HCl调pH至7.5
    \item 定容至1000 mL
\end{itemize}

\section{补充实验结果}

% 补充实验结果

\subsection{酶的动力学参数}

表 \ref{tab:kinetics} 展示了纯化后GPPS和LIS的动力学参数。

\begin{table}[h]
  \centering
  \bicaption{酶的动力学参数}{Kinetic parameters of enzymes}
  \label{tab:kinetics}
  \begin{tabular}{lccc}
    \toprule
    \textbf{酶} & \textbf{$K_m$ (mM)} & \textbf{$k_{cat}$ (s$^{-1}$)} & \textbf{$k_{cat}/K_m$ (M$^{-1}$s$^{-1}$)} \\
    \midrule
    GPPS & 0.15 & 12.5 & 8.33 $\times$ 10$^4$ \\
    LIS & 0.28 & 8.6 & 3.07 $\times$ 10$^4$ \\
    \bottomrule
  \end{tabular}
\end{table}

\subsection{固定化酶的储存稳定性}

固定化酶在4°C条件下储存30天后，GPPS保持82.5\%的初始活性，LIS保持78.3\%的初始活性。
 \clearpage

%%%%%%%%%%%%%%%%%致谢
\acknowledgement

本科学习的四年，是我人生中极为重要的阶段。值此毕业论文完成之际，我怀着无比感激的心情，向所有在学习和生活中给予我关怀、指导和帮助的人致以诚挚的谢意。

首先，衷心感谢我的论文指导老师——何新卫老师。老师不仅在选题、研究方法、论文结构等方面给予我极大的支持，更在整个写作过程中以严谨的治学态度、深入的专业素养和细致耐心的指导为我树立了榜样。从文献梳理到实验设计，从语言表达到逻辑推演，老师始终给予我耐心点拨，使我在科研能力、写作规范和独立思考方面都有了显著提升。

其次，感谢学院的各位任课老师，是你们系统、扎实的课程教学、是你们的悉心讲授和无私分享，使我在本科阶段打下了坚实的理论基础，拓展了专业视野，培养了严谨的思维方式和问题意识，也让我对人工智能领域产生了浓厚兴趣，为本课题的研究提供了有力支撑。

我还要特别感谢我的家人，是你们始终如一的理解、包容与支持，给予我稳定的后盾和精神的鼓励，让我能够安心学习，顺利完成本科学业和毕业论文。

感谢我的室友和朋友们，是你们让我的整个大学生活不只是学习，还有欢笑和陪伴。每一个深夜抱怨写不出论文、模型又跑崩了的时刻，是你们陪我“摆烂”也拉我重启；有时候，一个寝室灯都没关，大家“秉烛夜战”的样子，早已成了我最宝贵的青春记忆。

也感谢那些在路上遇到的陌生人，比如论坛上解答过我问题的网友，开源社区中那些无私共享的代码作者，是你们让我明白科研不仅是一种能力，也是一种传承与分享的精神。

这一切琐碎、鸡零狗碎又真切的经历，共同组成了我这一段难忘的毕业旅程。毕业并不是终点，而是另一个起点。感谢所有出现在我大学生活中的人和事，愿我们在各自的人生轨道上，都能勇敢、自由地奔跑。


\hfill Jun, 2025 

\hfill 张三














\end{document}
